\chapter{酸碱反应原理}

\section{酸碱理论}

\subsection{酸碱理论的发展脉络}
由浅入深、由低级到高级、由感性到理性，核心理论包括：
\begin{itemize}
    \item 阿累尼乌斯酸碱理论
    \item 酸碱质子理论  
    \item 酸碱电子理论
\end{itemize}

\subsection{阿累尼乌斯酸碱理论}
\subsubsection{主要观点}
在水溶液中解离时：
\begin{itemize}
    \item 酸（acid）：产生的阳离子\textbf{全部是} \ce{H+} 的物质
    \item 碱（base）：产生的阴离子\textbf{全部是} \ce{OH-} 的物质
\end{itemize}

\subsubsection{优缺点}
\begin{itemize}
    \item 优点：首次揭示酸碱的本质
    \item 局限性：
    \begin{enumerate}
        \item 割裂酸碱的关联性
        \item 无法解释非 \ce{H+}/\ce{OH-} 体系的酸碱性（如 \ce{NH4Cl} 溶液呈酸性、\ce{Na2CO3} 溶液呈碱性）
    \end{enumerate}
\end{itemize}

\subsection{酸碱质子理论}
\subsubsection{主要观点}
\begin{itemize}
    \item 酸：\textbf{能给出质子（\ce{H+}）}的物质
    \item 碱：\textbf{能接受质子（\ce{H+}）}的物质
    \item 两性物质：既能给出又能接受质子的物质
\end{itemize}

\subsubsection{共轭酸碱对}
酸给出1个质子后转变为\textbf{共轭碱}，碱接受1个质子后转变为\textbf{共轭酸}，二者构成共轭酸碱对。

\subsubsection{酸碱反应实质}
\textbf{两个共轭酸碱对之间的质子传递反应}

\subsubsection{适用范围扩展}
电解质解离、中和反应、盐的水解等本质均为质子传递反应。

\subsubsection{优缺点}
\begin{itemize}
    \item 优点：
    \begin{enumerate}
        \item 扩大酸碱范围（不再局限于水溶液）
        \item 明确反应本质是质子传递
    \end{enumerate}
    \item 局限性：仅适用于含质子的体系
\end{itemize}

\subsection{路易斯酸碱电子理论}

\begin{itemize}
    \item 酸：\textbf{接受电子对}的物质（电子对受体）
    \item 碱：\textbf{给出电子对}的物质（电子对给体）
    \item 反应实质：形成酸碱配合物
\end{itemize}


\section{单相解离平衡}

\subsection{水的解离平衡与溶液的pH值}
核心为水的离子积 $K_{\text{w}}^{\ominus} = c_{\text{r}}(\ce{H+}) \cdot c_{\text{r}}(\ce{OH-})$，$25^\circ \text{C}$ 时 $K_{\text{w}}^{\ominus} = 1.0 \times 10^{-14}$，$\text{pH} = -\lg c_{\text{r}}(\ce{H+})$。

\subsection{弱酸、弱碱的解离平衡}
\subsubsection{定义}
弱酸/弱碱与溶剂水的质子传递平衡，统称弱酸/弱碱的解离平衡。

\subsection{解离常数 \texorpdfstring{$K_{\text{i}}^{\ominus}$}{}}
对一元弱酸 \ce{HX <=> H+ + X-}：
$$K_{\text{i}}^{\ominus} = \frac{c_{\text{r}}(\ce{H+}) \cdot c_{\text{r}}(\ce{X-})}{c_{\text{r}}(\ce{HX})}$$

\subsubsection{解离度与稀释定律}
解离度 $\alpha$：弱电解质达平衡时\textbf{解离的百分数}：
$$\alpha = \frac{\text{已解离的弱电解质浓度}}{\text{初始浓度}} \times 100\%$$

与解离常数的关系：
$$K_{\text{i}}^{\ominus} = \frac{c\alpha^2}{1 - \alpha}$$
当 $\alpha < 5\%$ 时：
$$\alpha = \sqrt{\frac{K_{\text{i}}^{\ominus}}{c}}$$

\subsubsection{一元弱酸/弱碱溶液的离子浓度计算}
\begin{itemize}
    \item 一元弱酸：$c(\ce{H+}) = \sqrt{K_{\text{a}}^{\ominus} \cdot c}$
    \item 一元弱碱：$c(\ce{OH-}) = \sqrt{K_{\text{b}}^{\ominus} \cdot c}$
\end{itemize}

\subsubsection{多元弱酸的解离平衡}
\textbf{分步解离}，且一级解离常数 $\gg$ 二级解离常数 $\gg$ 三级解离常数。

\subsection{盐溶液的酸碱平衡（盐的水解）}
\subsubsection{定义}
盐的组成离子与水发生\textbf{质子传递反应}，使溶液呈现酸碱性。

\subsubsection{盐的分类及水解规律}
\begin{table}[ht]
\centering
\caption*{盐类水解类型与特征}
\begin{tabular}{lccc}
\toprule
盐的类型 & 示例 & 水解离子 & 溶液酸碱性 \\
\midrule
强酸强碱盐 & \ce{NaCl} & 无 & 中性 \\
强酸弱碱盐 & \ce{NH4Cl} & 阳离子 & 酸性 \\
弱酸强碱盐 & \ce{NaAc} & 阴离子 & 碱性 \\
弱酸弱碱盐 & \ce{NH4Ac} & 阴阳离子 & 视情况而定 \\
\bottomrule
\end{tabular}
\end{table}

\subsubsection{强酸弱碱盐的水解计算}
以 \ce{NH4Cl} 为例：
$$c(\ce{H+}) = \sqrt{\frac{K_{\text{w}}^{\ominus} \cdot c_{\text{盐}}}{K_{\text{b}}^{\ominus}(\ce{NH3})}}$$

\subsubsection{弱酸强碱盐的水解计算}
以 \ce{NaAc} 为例：
$$c(\ce{OH-}) = \sqrt{\frac{K_{\text{w}}^{\ominus} \cdot c_{\text{盐}}}{K_{\text{a}}^{\ominus}(\ce{HAc})}}$$

\subsubsection{弱酸弱碱盐的水解计算}
以 \ce{NH4Ac} 为例：
$$c(\ce{H+}) = \sqrt{\frac{K_{\text{w}}^{\ominus} \cdot K_{\text{a}}^{\ominus}(\ce{HAc})}{K_{\text{b}}^{\ominus}(\ce{NH3})}}$$

\subsubsection{影响盐类水解的因素}
\begin{enumerate}
    \item 盐的本性：组成盐的酸/碱越弱，水解程度越大
    \item 溶液酸碱性：加酸/碱可抑制相应盐的水解
    \item 温度：水解吸热，升温促进水解
\end{enumerate}

\subsection{同离子效应与缓冲溶液}
\subsubsection{同离子效应}
在弱电解质溶液中加入\textbf{含相同离子的强电解质}，使弱电解质的解离度降低的现象。

\subsubsection{缓冲溶液的定义}
能\textbf{抵抗外加少量强酸、强碱或稀释}，保持自身pH基本不变的溶液。

\subsubsection{缓冲对的组成}
由\textbf{弱酸-共轭碱}或\textbf{弱碱-共轭酸}组成。

\subsubsection{缓冲原理}
以 \ce{HAc-NaAc} 为例：
\begin{itemize}
    \item 加少量强酸：\ce{Ac- + H+ <=> HAc}
    \item 加少量强碱：\ce{HAc + OH- <=> Ac- + H2O}
\end{itemize}

\subsubsection{缓冲溶液pH计算}
\begin{itemize}
    \item 弱酸-共轭碱体系：$\text{pH} = \text{p}K_{\text{a}}^{\ominus} - \lg \frac{c_{\text{酸}}}{c_{\text{盐}}}$
    \item 弱碱-共轭酸体系：$\text{pOH} = \text{p}K_{\text{b}}^{\ominus} - \lg \frac{c_{\text{碱}}}{c_{\text{盐}}}$
\end{itemize}

\subsubsection{缓冲溶液的关键性质}
\begin{enumerate}
    \item 缓冲能力有限
    \item 当 $\dfrac{c_{\text{酸}}}{c_{\text{盐}}} = 1$ 时缓冲能力最强
    \item 选择 $\text{p}K_{\text{a}}^{\ominus}$ 与目标pH接近的缓冲对
\end{enumerate}

\subsubsection{酸碱指示剂}
弱有机酸/碱，变色范围：$\text{pH} = \text{p}K_{\text{a}}^{\ominus} \pm 1$

\subsection{强电解质的解离与表观解离度}
\subsubsection{矛盾点}
强电解质理论上应$100\%$解离，但实测解离度$<100\%$。

\subsubsection{德拜-休克尔理论解释}
强电解质溶液中存在\textbf{离子氛}，离子移动受阻。

\subsubsection{影响离子氛的因素}
\begin{enumerate}
    \item 浓度：浓度越大，阻碍越强
    \item 离子电荷：电荷越高，阻碍越强
\end{enumerate}

\section{配位化合物与配位平衡}

\subsection{配位化学概述}
\subsubsection{奠基人}
维尔纳（A. Werner），1893年提出配位理论。

\subsubsection{重要应用}
\begin{itemize}
    \item 催化反应
    \item 生物过程：血红蛋白的氧合与携氧
    \item 化学制备：如银氨配离子的生成
\end{itemize}

\subsection{配位化合物的基本概念}
\subsubsection{定义}
由\textbf{配体}与\textbf{中心原子/离子}按一定组成和空间构型形成的化合物。

\subsubsection{组成}
\begin{itemize}
    \item 中心离子/原子：接受电子对的核心
    \item 配体：围绕中心离子/原子的中性分子或阴离子
    \item 配位原子：配体中直接提供孤对电子的原子
    \item 配位数：与中心离子/原子直接成键的配位原子个数
    \item 内界与外界
\end{itemize}

\subsubsection{配体的分类}
\begin{itemize}
    \item 单齿配体：1个配体含1个配位原子
    \item 多齿配体：1个配体含2个及以上配位原子
\end{itemize}

\subsubsection{配位数的特点与影响因素}
\begin{itemize}
    \item 常见配位数：4和6
    \item 影响因素：中心离子/原子的电荷与半径、配体的电荷与半径、温度与配体浓度
\end{itemize}

\subsection{配位化合物的命名}
\subsubsection{基本原则}
\begin{itemize}
    \item 内界与外界：负离子在前，正离子在后
    \item 内界：配体合中心离子，配体前标数目，中心离子后标氧化数
\end{itemize}

\subsubsection{内界多配体的命名顺序}
\begin{enumerate}
    \item 先无机配体，后有机配体
    \item 先阴离子配体，后阳离子配体、分子配体
    \item 同类配体按配位原子元素符号英文字母顺序
\end{enumerate}

\subsection{配离子解离平衡}
\subsubsection{解离平衡与平衡常数}
以 \ce{[Cu(NH3)4]^{2+}} 为例：
\begin{itemize}
    \item 不稳定常数：$K_{\text{d}}^{\ominus} = \dfrac{c_{\text{r}}(\ce{Cu^{2+}}) \cdot c_{\text{r}}^4(\ce{NH3})}{c_{\text{r}}(\ce{[Cu(NH3)4]^{2+}})}$
    \item 稳定常数：$K_{\text{f}}^{\ominus} = \dfrac{1}{K_{\text{d}}^{\ominus}}$
\end{itemize}

\subsubsection{逐级稳定常数}
配离子生成是\textbf{分步进行}的，总稳定常数 $K_{\text{f}}^{\ominus} = \prod _1^n \beta_n$。

\subsubsection{稳定常数的应用}
比较同类型配合物的稳定性。

\subsubsection{配离子解离平衡的移动}
\begin{enumerate}
    \item 与难溶盐沉淀的相互作用
    \item 生成更稳定的配离子
    \item 氧化还原反应的影响
\end{enumerate}